\documentclass[11pt,a4paper]{article}
\usepackage[margin=2.2cm]{geometry}
\usepackage{fontspec}
\usepackage[spanish,es-nodecimaldot]{babel}
\usepackage{xcolor}
\usepackage{titlesec}
\usepackage{enumitem}
\usepackage{tabularx}
\usepackage{booktabs}
\usepackage{fancyhdr}
\usepackage{hyperref}
\usepackage{microtype}
\usepackage{tcolorbox}
\usepackage{lastpage}

\definecolor{navy}{HTML}{17324D}
\definecolor{blue}{HTML}{2878A8}
\definecolor{lightblue}{HTML}{EAF3F8}
\definecolor{green}{HTML}{24735C}
\definecolor{amber}{HTML}{B56A16}
\definecolor{lightgray}{HTML}{F3F5F7}
\definecolor{midgray}{HTML}{66727D}
\hypersetup{colorlinks=true,linkcolor=blue,urlcolor=blue}
\setlength{\parindent}{0pt}
\setlength{\parskip}{6pt}
\renewcommand{\arraystretch}{1.12}
\setlist[itemize]{leftmargin=1.4em,itemsep=2pt,topsep=3pt}
\setlist[enumerate]{leftmargin=1.6em,itemsep=3pt,topsep=3pt}
\titleformat{\section}{\Large\bfseries\color{navy}}{}{0pt}{}[\vspace{-3pt}\color{blue}\rule{\linewidth}{0.6pt}]
\titleformat{\subsection}{\large\bfseries\color{blue}}{}{0pt}{}
\pagestyle{fancy}
\fancyhf{}
\lhead{\small\color{midgray}LatAm Macro Monitor}
\rhead{\small\color{midgray}Estado del proyecto}
\cfoot{\small\color{midgray}Página \thepage\ de \pageref{LastPage}}
\renewcommand{\headrulewidth}{0.3pt}

\newtcolorbox{summarybox}{colback=lightblue,colframe=blue,boxrule=0.7pt,arc=2mm,left=4mm,right=4mm,top=3mm,bottom=3mm}
\newtcolorbox{warningbox}{colback=amber!7,colframe=amber,boxrule=0.7pt,arc=2mm,left=4mm,right=4mm,top=3mm,bottom=3mm}
\newcommand{\code}[1]{\texttt{#1}}
\newcommand{\statusdone}{\textcolor{green}{\textbf{Hecho}}}
\newcommand{\statuspartial}{\textcolor{amber}{\textbf{Parcial}}}
\newcommand{\statuspending}{\textcolor{midgray}{\textbf{Pendiente}}}

\begin{document}

\begin{titlepage}
\pagecolor{navy}
\color{white}
\vspace*{2.2cm}
{\Large INFORME DE AVANCE\par}
\vspace{0.5cm}
{\fontsize{30}{36}\selectfont\bfseries LatAm Macro Monitor\par}
\vspace{0.4cm}
{\Large Explicación completa del proyecto Economics\par}
\vspace{1.4cm}
\color{white!85}
{\large Qué se construyó, cómo funciona y qué falta\par}
\vfill
\begin{tcolorbox}[colback=white,colframe=white!35,coltext=navy,boxrule=0.5pt,arc=2mm]
\textbf{Estado observado:} arquitectura inicial en desarrollo\\
\textbf{Rama actual:} \code{features/project.config}\\
\textbf{Fecha de revisión:} 20 de julio de 2026\\
\textbf{Base del análisis:} archivos actuales, historial Git y cambio local sin confirmar
\end{tcolorbox}
\vspace{1cm}
{\small Documento generado a partir del contenido real del repositorio.}
\end{titlepage}
\nopagecolor
\color{black}

\section{Resumen ejecutivo}

\begin{summarybox}
El proyecto define la base de una plataforma abierta para monitorear indicadores macroeconómicos y financieros de América Latina. Hasta ahora se ha construido el esqueleto del repositorio, la configuración central, el catálogo geográfico y una interfaz común para futuras fuentes de datos. Aún no existe un flujo funcional que descargue, procese, analice o publique datos.
\end{summarybox}

La intención del proyecto está expresada en el \code{README.md}: recolectar datos macroeconómicos públicos, monitorear mercados financieros latinoamericanos, producir reportes reproducibles con Quarto y publicar un sitio automatizado mediante GitHub Pages.

El avance real puede resumirse así:

\begin{itemize}
  \item \statusdone: definición del propósito, estructura de carpetas, exclusiones de Git, configuración de rutas y lista inicial de países.
  \item \statuspartial: arquitectura para proveedores de datos y conector FRED, todavía sin implementación.
  \item \statuspending: descarga efectiva, transformación, indicadores, gráficos, notebooks, pruebas, reportes Quarto, automatización y publicación web.
\end{itemize}

\subsection{Lectura rápida del estado}

\begin{tabularx}{\textwidth}{@{}p{3.5cm}X p{2.2cm}@{}}
\toprule
\textbf{Área} & \textbf{Evidencia encontrada} & \textbf{Estado} \\
\midrule
Propósito & README con nombre, descripción y cuatro objetivos & \statusdone \\
Configuración & Rutas, fecha inicial y nombre del proyecto & \statusdone \\
Cobertura geográfica & 11 economías definidas & \statusdone \\
Fuentes de datos & Clase abstracta y archivos para FRED, OECD, World Bank y Yahoo & \statuspartial \\
Datos y análisis & Carpetas creadas, pero sin contenido rastreado & \statuspending \\
Calidad y entrega & Sin pruebas, reportes, CI ni sitio publicado & \statuspending \\
\bottomrule
\end{tabularx}

\section{Qué se ha hecho, en orden cronológico}

El historial contiene tres commits. Todos fueron realizados el 6 de julio de 2026 y muestran una progresión lógica desde la estructura general hacia la capa de datos.

\subsection{1. Inicialización del proyecto}

Commit \code{7c1574b}: \emph{Initialize project structure}.

\begin{itemize}
  \item Se creó el \code{README.md} con el nombre \textbf{LatAm Macro Monitor}, una descripción breve, el estado ``Under development'' y los objetivos principales.
  \item Se agregó \code{.gitignore} para excluir el entorno virtual, secretos de \code{.env}, cachés de Python, checkpoints de Jupyter, configuración de VS Code y datos crudos/procesados.
  \item Se creó \code{proyect.toml}, todavía vacío. El nombre también parece provisional: en el ecosistema Python normalmente sería \code{pyproject.toml}.
  \item En el sistema de archivos quedaron preparadas las carpetas \code{data}, \code{notebooks}, \code{reports}, \code{src} y \code{test}.
\end{itemize}

\textbf{Qué aporta:} una separación inicial entre código, datos, análisis, reportes y pruebas, además de evitar que archivos sensibles o pesados entren accidentalmente en Git.

\subsection{2. Módulo de configuración}

Commit \code{2f2496a}: \emph{add config module}.

Se creó el paquete \code{src/config} con dos piezas:

\begin{itemize}
  \item \code{settings.py} calcula la raíz del repositorio desde la ubicación del archivo y deriva rutas para \code{data}, \code{data/raw}, \code{data/processed} y \code{reports}. También guarda \code{Project\_name = "Economics"} y \code{Start\_date = "2026-07-06"}.
  \item \code{countries.py} define un diccionario de 11 economías: Chile, Brasil, México, Colombia, Perú, Argentina, Estados Unidos, Eurozona, China, Corea del Sur y Japón.
\end{itemize}

Esta selección combina mercados latinoamericanos con economías externas relevantes como referencia, socios comerciales o impulsores globales.

\subsection{3. Arquitectura base de fuentes de datos}

Commit \code{9457588}: \emph{define base archiqueture for data sources}.

\begin{itemize}
  \item \code{src/data/base.py} introduce la clase abstracta \code{datasource}.
  \item Su método abstracto \code{download\_data(...)} establece el contrato: cada proveedor deberá descargar información y devolver un \code{pandas.DataFrame}.
  \item Se crearon módulos reservados para \code{fred.py}, \code{oecd.py}, \code{word\_bank.py} y \code{yahoo.py}.
  \item \code{fred.py} contiene una subclase, pero su método todavía lanza \code{NotImplementedError}; por lo tanto, no descarga nada aún.
\end{itemize}

\textbf{Qué aporta:} una futura API uniforme. El resto del proyecto podría pedir datos a distintos proveedores mediante la misma operación, reduciendo acoplamiento y facilitando pruebas o sustituciones.

\section{Cómo está organizado el proyecto}

\begin{tabularx}{\textwidth}{@{}p{4.0cm}X@{}}
\toprule
\textbf{Ruta} & \textbf{Función actual o prevista} \\
\midrule
\code{README.md} & Presenta la visión y los objetivos. \\
\code{.gitignore} & Protege secretos, entorno local, cachés y datos generados. \\
\code{data/raw} & Destino previsto para datos originales; actualmente vacío e ignorado. \\
\code{data/processed} & Destino previsto para datos transformados; actualmente vacío e ignorado. \\
\code{notebooks} & Espacio previsto para exploración; actualmente vacío. \\
\code{reports} & Espacio previsto para reportes reproducibles; actualmente vacío. \\
\code{src/config} & Configuración central y catálogo de países; ya contiene código. \\
\code{src/data} & Contrato de proveedores y conectores previstos; implementación parcial. \\
\code{src/indicators} & Futuro cálculo de indicadores; actualmente vacío. \\
\code{src/processing} & Futuras limpiezas y transformaciones; actualmente vacío. \\
\code{src/visualizacion} & Futuras visualizaciones; actualmente vacío. \\
\code{src/utils} & Futuras utilidades compartidas; actualmente vacío. \\
\code{test} & Futuras pruebas automatizadas; actualmente vacío. \\
\bottomrule
\end{tabularx}

\subsection{Flujo previsto}

La arquitectura sugiere este recorrido, aunque todavía no está implementado:

\begin{center}
\fcolorbox{blue}{lightblue}{\strut Fuentes públicas}
\quad$\rightarrow$\quad
\fcolorbox{blue}{lightblue}{\strut \code{src/data}}
\quad$\rightarrow$\quad
\fcolorbox{blue}{lightblue}{\strut \code{data/raw}}
\quad$\rightarrow$\quad
\fcolorbox{blue}{lightblue}{\strut procesamiento}
\end{center}
\begin{center}
$\downarrow$\\[2pt]
\fcolorbox{green}{green!7}{\strut indicadores y visualizaciones}
\quad$\rightarrow$\quad
\fcolorbox{green}{green!7}{\strut reportes Quarto}
\quad$\rightarrow$\quad
\fcolorbox{green}{green!7}{\strut GitHub Pages}
\end{center}

\section{Decisiones técnicas ya tomadas}

\subsection{Python y pandas como base}

La anotación de retorno de \code{download\_data} fija \code{pandas.DataFrame} como formato común. Esto es adecuado para series de tiempo tabulares y permite que las capas posteriores operen de forma homogénea.

\subsection{Interfaz abstracta para proveedores}

El uso de \code{ABC} y \code{@abstractmethod} impide conceptualizar cada fuente como un script aislado. La decisión apunta a conectores intercambiables y a una arquitectura extensible.

\subsection{Configuración centralizada}

Las rutas se derivan con \code{pathlib.Path}, evitando rutas absolutas ligadas a un solo computador. Esto favorece la reproducibilidad, aunque los nombres de variables deberían normalizarse.

\subsection{Separación de datos crudos y procesados}

La división \code{raw/processed} preserva el insumo original y reserva otro espacio para resultados limpios. Al ignorar ambas carpetas, Git no almacenará automáticamente datos locales; más adelante habrá que definir cómo se reproducen o distribuyen.

\subsection{Git y trabajo por rama}

La rama actual es \code{features/project.config}. El commit más reciente está localmente por delante de su referencia remota, mientras que \code{src/data/fred.py} tiene una modificación adicional sin commit.

\section{Hallazgos y riesgos actuales}

\begin{warningbox}
\textbf{El proyecto todavía no es ejecutable de extremo a extremo.} Su valor actual es estructural: expresa una intención y prepara componentes, pero no produce todavía datos, gráficos, reportes ni una web.
\end{warningbox}

\begin{enumerate}
  \item \textbf{Importación local problemática en FRED.} El commit usaba una importación relativa desde \code{.base}, apropiada al importar el módulo como paquete. El cambio local la reemplaza por una importación desde \code{base}; al cargar \code{src.data.fred}, Python normalmente no encuentra ese módulo superior. Conviene restaurar la importación relativa o adoptar importaciones absolutas coherentes.
  \item \textbf{Conectores vacíos.} OECD, World Bank y Yahoo son archivos de cero bytes. FRED solo contiene un marcador que lanza una excepción.
  \item \textbf{Configuración de paquete incompleta.} \code{proyect.toml} está vacío y no declara dependencias, versión, empaquetado ni herramientas. Esto impide recrear el entorno de manera confiable.
  \item \textbf{Convenciones inconsistentes.} Hay códigos de país de dos y tres letras mezclados (por ejemplo, \code{CL} y \code{BRA}); \code{word\_bank.py} probablemente debería llamarse \code{world\_bank.py}; y variables/clases usan mayúsculas de forma irregular.
  \item \textbf{Sin pruebas.} No hay comprobaciones para rutas, catálogo de países, contrato de proveedores ni esquemas de datos.
  \item \textbf{Sin persistencia reproducible.} Las carpetas de datos están ignoradas y no hay scripts, manifiestos o instrucciones para regenerarlas.
  \item \textbf{README breve.} No explica instalación, ejecución, fuentes, variables de entorno, estructura de datos ni flujo de contribución.
  \item \textbf{Entorno virtual no verificable en esta revisión.} El ejecutable registrado por \code{.venv} apunta a una instalación de Python 3.14 que no estuvo accesible desde el entorno de revisión; esto puede ser una limitación del entorno o un entorno virtual roto.
\end{enumerate}

\section{Qué falta para alcanzar los objetivos}

\begin{tabularx}{\textwidth}{@{}p{4.2cm}X@{}}
\toprule
\textbf{Objetivo declarado} & \textbf{Trabajo restante principal} \\
\midrule
Recolectar datos públicos & Elegir series, credenciales cuando correspondan, implementar conectores, normalizar fechas/unidades, manejar errores y guardar resultados. \\
Monitorear mercados & Definir indicadores, frecuencias, monedas, benchmarks y reglas de actualización; implementar cálculos y validaciones. \\
Reportes con Quarto & Crear archivos \code{.qmd}, tablas y gráficos; parametrizar fechas/países y documentar cómo renderizarlos. \\
Sitio con GitHub Pages & Configurar Quarto website, workflow de GitHub Actions, caché/datos, publicación y manejo seguro de secretos. \\
\bottomrule
\end{tabularx}

\subsection{Secuencia recomendada}

\begin{enumerate}
  \item \textbf{Hacer instalable el proyecto:} renombrar y completar \code{pyproject.toml}, declarar dependencias y documentar la instalación.
  \item \textbf{Normalizar fundamentos:} convenciones de nombres, códigos geográficos, imports y estructura de paquetes.
  \item \textbf{Entregar un flujo vertical mínimo:} seleccionar una serie FRED, descargarla, validarla, guardarla y mostrarla en un reporte sencillo.
  \item \textbf{Añadir pruebas:} rutas, conectores, columnas esperadas, frecuencia, ausencia de duplicados y manejo de fallas de red.
  \item \textbf{Extender proveedores:} World Bank, OECD y Yahoo con el mismo contrato y un esquema común.
  \item \textbf{Construir indicadores y visualizaciones:} primero para uno o dos países, luego ampliar cobertura.
  \item \textbf{Automatizar:} actualización programada, Quarto y GitHub Pages cuando el flujo local sea estable.
\end{enumerate}

\section{Conclusión}

Has hecho correctamente el trabajo de cimentación: definiste qué producto quieres construir, separaste responsabilidades, centralizaste rutas, delimitaste la cobertura geográfica e introdujiste una abstracción común para fuentes de datos. Los tres commits muestran una evolución ordenada y comprensible.

El proyecto aún está antes de su primer producto mínimo funcional. El siguiente hito más útil no sería crear más carpetas, sino completar una cadena pequeña y verificable desde una fuente real hasta un reporte. Esa primera cadena permitirá validar las decisiones de arquitectura antes de replicarlas en todos los países y proveedores.

\vfill
{\small\color{midgray}\textbf{Nota metodológica.} Este informe describe el estado observado del repositorio al 20 de julio de 2026. No atribuye como terminadas capacidades que solo aparecen como objetivos, carpetas vacías o marcadores de implementación. El PDF es el único artefacto nuevo de entrega; no se modificó el código fuente del proyecto.}

\end{document}
